



 \chapter{四福音背景}
\section{歷史背景}
\subsection{引用舊約聖經}
\subsubsection{馬太福音}
% Table generated by Excel2LaTeX from sheet 't1'
\begin{table}[htbp]
  \centering 
    %\begin{tabular}{|p{10.785em}|p{28.855em}|p{6.93em}|}
   % \toprule
    %\rowcolor[rgb]{ .749,  .749,  .749} \textcolor[rgb]{ .267,  .447,  .769}{\textbf{馬太福音}} 
    %& \textcolor[rgb]{ .267,  .447,  .769}{\textbf{馬太引用經文內容}} 
    %& \textcolor[rgb]{ .267,  .447,  .769}{\textbf{引用出處}} \\
    %\midrule
    \begin{tabular}{|p{4.2em}|p{31.645em}|p{7.3em}|}
    \toprule
    \rowcolor[rgb]{ .749,  .749,  .749} \textcolor[rgb]{ .267,  .447,  .769}{\textbf{馬太福音}} 
    & \textcolor[rgb]{ .267,  .447,  .769}{\textbf{馬太引用經文內容}} 
    & \textcolor[rgb]{ .267,  .447,  .769}{\textbf{引用出處}} \\    
  %\begin{tabular}{|p{4.07em}|p{31.07em}|p{7.355em}|}
  %  \toprule
  %  \rowcolor[rgb]{ .749,  .749,  .749} \textcolor[rgb]{ .267,  .447,  .769}{\textbf{馬太福音}}
   %  & \textcolor[rgb]{ .267,  .447,  .769}{\textbf{馬太引用經文內容}} 
   %  & \textcolor[rgb]{ .267,  .447,  .769}{\textbf{引用出處}} \\
    \midrule        
    1:1-6 & 家譜    & 代上 2:1-14，3:5-24 \\
    \midrule
    1:22  & 必有童女懷孕生子；人要稱他的名為以馬內利 & 以賽亞書7：14 \\
    \midrule
    2:6   & 猶大地的伯利恆啊，你在猶大諸城中並不是最小的；因為將來有一位君王要從你那裡出來，牧養我以色列民。 & 彌迦書 5:2 \\
    \midrule
    \multirow{3}[6]{*}{2:15} & 住在那裡，直到希律死了。這是要應驗主藉先知所說的話說： & 出埃及記 4:22 \\
\cmidrule{3-3}    \multicolumn{1}{|l|}{} & \multirow{2}[4]{*}{「我從埃及召出我的兒子來。」} & 民數記 24:8 \\
\cmidrule{3-3}    \multicolumn{1}{|l|}{} & \multicolumn{1}{l|}{} & 何西阿書 11:1 \\
    \midrule
    2:18  & 在拉瑪聽見號咷大哭的聲音，是拉結哭他兒女，不肯受安慰，因為他們都不在了。 & 耶利米書 31:15 \\
    \midrule
    2:23  & 這是要應驗先知所說，他將稱為拿撒勒人的話了 & ? \\
    \midrule
    3:3   & 在曠野有人聲喊著說：預備主的道，修直他的路！ & 以賽亞書 40:3 \\
    \midrule
    4:4   & 人活著，不是單靠食物，乃是靠神口裡所出的一切話。 & 申命記 8:3 \\
    \midrule
    4:6   & 主要為你吩咐他的使者用手托著你，免得你的腳碰在石頭上。 & 詩篇 91:11，12 \\
    \midrule
    4:7   & 不可試探主你的神。 & 申命記 6:16 \\
    \midrule
    4:10  & 當拜主你的神，單要事奉他 & 申命記 6:13， 10:20 \\
    \midrule
    4:15  & 西布倫地，拿弗他利地，就是沿海的路，約但河外，外邦人的加利利地，那坐在黑暗裡的百姓看見了大光；坐在死蔭之地的人有光發現照著他們。 & 以賽亞書 9:1 \\
    \bottomrule
    \end{tabular}%
  \label{tab:addlabel}%
    \caption{馬太福音引用舊約表格 1}
\end{table}%





% Table generated by Excel2LaTeX from sheet 't2'
\begin{table}[htbp]
  \centering
  %  \begin{tabular}{|p{4.07em}|p{31.07em}|p{7.355em}|}
  %  \toprule
  %  \rowcolor[rgb]{ .749,  .749,  .749} \textcolor[rgb]{ .267,  .447,  .769}{\textbf{馬太福音}}
   %  & \textcolor[rgb]{ .267,  .447,  .769}{\textbf{馬太引用經文內容}} 
  %   & \textcolor[rgb]{ .267,  .447,  .769}{\textbf{引用出處}} \\
  %  \midrule
    \begin{tabular}{|p{4.2em}|p{31.645em}|p{7.3em}|}
    \toprule
    \rowcolor[rgb]{ .749,  .749,  .749} \textcolor[rgb]{ .267,  .447,  .769}{\textbf{馬太福音}} 
    & \textcolor[rgb]{ .267,  .447,  .769}{\textbf{馬太引用經文內容}} 
    & \textcolor[rgb]{ .267,  .447,  .769}{\textbf{引用出處}} \\    
  \midrule    
    \multirow{2}[4]{*}{5:21} & \multirow{2}[4]{*}{不可殺人,凡殺人的難免受審判。} & 出埃及記 20:13 \\
\cmidrule{3-3}    \multicolumn{1}{|l|}{} & \multicolumn{1}{l|}{} & 申命記 5:17 \\
    \midrule
    \multirow{3}[6]{*}{5:27} & \multirow{3}[6]{*}{不可姦淫。} & 申命記 5:18 \\
\cmidrule{3-3}    \multicolumn{1}{|l|}{} & \multicolumn{1}{l|}{} & 利未記 18:20 \\
\cmidrule{3-3}    \multicolumn{1}{|l|}{} & \multicolumn{1}{l|}{} & 出埃及記 20:14 \\
    \midrule
    5:31  & 人若休妻，就當給他休書。 & 申命記 24:1，3 \\
    \midrule
    \multicolumn{1}{|l|}{\multirow{2}[4]{*}{5:33}} & \multirow{2}[4]{*}{不可背誓，所起的誓總要向主謹守。} & 民數記 30:2 \\
\cmidrule{3-3}    \multicolumn{1}{|l|}{} & \multicolumn{1}{l|}{} & 申命記 23:23 \\
    \midrule
    \multicolumn{1}{|l|}{\multirow{3}[6]{*}{5:38}} & \multirow{3}[6]{*}{以眼還眼，以牙還牙。} & 利未記 24:20 \\
\cmidrule{3-3}    \multicolumn{1}{|l|}{} & \multicolumn{1}{l|}{} & 申命記 19:21 \\
\cmidrule{3-3}    \multicolumn{1}{|l|}{} & \multicolumn{1}{l|}{} & 出埃及記 21:24 \\
    \midrule
    \multicolumn{1}{|l|}{5:43} & 當愛你的鄰舍，恨你的仇敵。 & 利未記 19:18 \\
        \midrule
        8:4   & 耶稣对他说：“你切不可告诉人，只要去把身体给祭司察看，献上摩西所吩咐的礼物，对众人做证据。” & 利未記 14:3,4 \\    
    \midrule
    8:17  & 他代替我們的軟弱，擔當我們的疾病。 & 以賽亞書 53:4 \\
    \midrule
    9:13  & 我喜愛憐恤，不喜愛祭祀。 & 何西阿書 6:6 \\
    \midrule
    11:5  & 瞎子看見，瘸子行走,長大痲瘋的潔淨，聾子聽見，死人復活，窮人有福音傳給他們。 & 以賽亞書 35：5 \\
    \midrule
    11:10 & 我要差遣我的使者在你前面預備道路 & 瑪拉基書 3:1 \\
    \midrule
    11:14 & 你們若肯領受，這人就是那應當來的以利亞。 & 瑪拉基書 3:1, 4:5 \\
    \midrule
    11:24 & 當審判的日子，所多瑪所受的，比你還容易受呢！ & 創世記19：1-29 \\
    \midrule
    12:2  & 看哪，你的門徒做安息日不可做的事了！ & 出埃及記20：8-11 \\
    \midrule
    \multirow{2}[4]{*}{12:3,4} & 經上記著大衛和跟從他的人飢餓之時所做的事，你們沒有念過嗎？他怎 & 利未記 24:9 \\
\cmidrule{3-3}    \multicolumn{1}{|l|}{} & 麼進了神的殿，吃了陳設餅，這餅不是他和跟從他的人可以吃得，惟獨祭司才可以吃。 & 撒母耳記上 21:4,6 \\
    \midrule
    12:5  & 當安息日，祭司在殿裡犯了安息日還是沒有罪，你們沒有念過嗎？ & 利未記 24:8 \\
    \midrule
    12:18-21 & 看哪！我的僕人，我所揀選，所親愛，心裡所喜悅的，我要將我的靈賜給他；他必將公理傳給外邦。他不爭競，不喧嚷；街上也沒有人聽見他的聲音。壓傷的蘆葦，他不折斷；將殘的燈火，他不吹滅；等他施行公理，叫公理得勝。外邦人都要仰望他的名。 & 以賽亞書 42:1-4 \\
    \bottomrule
    \end{tabular}%
  \label{tab:addlabel}%
   \caption{馬太福音引用舊約表格 2}
\end{table}%





% Table generated by Excel2LaTeX from sheet 't3'
\begin{table}[htbp]
  \centering
   \begin{tabular}{|p{4.93em}|p{31.em}|p{7. em}|}
    \toprule
    \rowcolor[rgb]{ .749,  .749,  .749} \textcolor[rgb]{ .267,  .447,  .769}{\textbf{馬太福音}} 
    & \textcolor[rgb]{ .267,  .447,  .769}{\textbf{馬太引用經文內容}} 
    & \textcolor[rgb]{ .267,  .447,  .769}{\textbf{引用出處}} \\
    \midrule
    \multicolumn{1}{|p{4.93em}|}{12:39-40} & 一個邪惡淫亂的世代求看神蹟，除了先知約拿的神蹟以外，再沒有神蹟給他們看。 約拿三日三夜在大魚肚腹中，人子也要這樣三日三夜在地裡頭。當審判的時候，尼尼微人要起來定這世代的罪，因為尼尼微人聽了約拿所傳的就悔改了。看哪，在這裡有一人比約拿更大！  & 約拿書  \\
    \midrule
    \multicolumn{1}{|l|}{\multirow{2}[4]{*}{12:42}} & 當審判的時候，南方的女王要起來定這世代的罪，因為他從地極而來， & 歷代志下 9:1 \\
\cmidrule{3-3}          & 要聽所羅門的智慧話。看哪！在這裡有一人比所羅門更大。 & 列王紀上 10:1 \\
    \midrule
    \multicolumn{1}{|p{4.93em}|}{13:15} & 因為這百姓油蒙了心，耳朵發沉，眼睛閉著，恐怕眼睛看見，耳朵聽見，心裡明白，回轉過來，我就醫治他們。 & 以賽亞書 6:10 \\
    \midrule
    \multicolumn{1}{|p{4.93em}|}{13:35} & 我要开口，用比喻说出创世以来被隐藏的事。 & 詩篇 78:2 \\
    \midrule
    \multicolumn{1}{|p{4.93em}|}{15:4} & 『當孝敬父母』，又說：『咒罵父母的，必治死他。 & 出埃及記 21:17 \\
    \midrule
    \multicolumn{1}{|p{4.93em}|}{15:8,9} & 『這百姓用嘴唇尊敬我，心卻遠離我。 9他們將人的吩咐當做道理教導人，所以拜我也是枉然。』 & 以賽亞書 29:13 \\
    \midrule
    \multicolumn{1}{|p{4.93em}|}{17:24} & 到了迦百農，有收丁稅的人來見彼得，說：「你們的先生不納丁稅嗎？」 & 出埃及記 30:13, 38:26 \\
    \midrule
    \multicolumn{1}{|p{4.93em}|}{19:4} & 那起初造人的，是‘造男造女 & 創世記 1:27,5:2 \\
    \midrule
    \multicolumn{1}{|l|}{\multirow{2}[4]{*}{19:5}} & \multirow{2}[4]{*}{因此，人要离开父母，与妻子联合，二人成为一体。} & 歷代志下 9:1 \\
\cmidrule{3-3}          & \multicolumn{1}{l|}{} & 瑪拉基書 2:15 \\
    \midrule
    \multicolumn{1}{|p{4.93em}|}{19:7} & 法利赛人说：“这样，摩西为什么吩咐给妻子休书，就可以休她呢？” & 申命記 24:1 \\
    \midrule
    \multicolumn{1}{|p{4.93em}|}{19:19} & 不 可 偷 盜 。 & 申命記 5:19 \\
    \midrule
    \multicolumn{1}{|l|}{\multirow{2}[4]{*}{19:19}} & \multirow{2}[4]{*}{不可作假見證} 
    	& 出埃及記 20:16 \\
\cmidrule{3-3}          & \multicolumn{1}{l|}{} &申命記 5:20 \\
%\cellcolor[rgb]{ .992,  .996,  1}申命記 5:20 \\

    \midrule
    \multicolumn{1}{|l|}{\multirow{2}[4]{*}{21:5}} & 要對錫安的居民說：『看哪，你的王來到你這裡，是溫柔的， &以賽亞書 62:11 \\
\cmidrule{3-3}          & 又騎著驢，就是騎著驢駒子。』 & 撒迦利亞書 9:9 \\
    \midrule
    \multicolumn{1}{|p{4.93em}|}{21:9} & 和散那歸於大衛的子孫！奉主名來的是應當稱頌的！高高在上和散那！ & 詩篇 118:26 \\
    \midrule
    \multicolumn{1}{|l|}{\multirow{2}[4]{*}{21:13}} & \multirow{2}[4]{*}{『我的殿必稱為禱告的殿』，你們倒使它成為賊窩了！} & 以賽亞書 56:7 \\
\cmidrule{3-3}          & \multicolumn{1}{l|}{} & 耶利米書 7:11 \\
    \midrule
    \multicolumn{1}{|p{4.93em}|}{21:16} & 『你從嬰孩和吃奶的口中完全了讚美 & 詩篇 8:2 \\
    \midrule
    \multicolumn{1}{|l|}{\multirow{2}[4]{*}{21:42}} & 『匠人所棄的石頭，已做了房角的頭塊石頭。這是主所做的， & 詩篇 118:22 \\
\cmidrule{3-3}          & 在我們眼中看為稀奇。』 & 以賽亞書 28:16 \\
    \midrule
    \multicolumn{1}{|l|}{\multirow{2}[4]{*}{22:32}} & 『我是亞伯拉罕的神、以撒的神、雅各的神。』 & 出埃及記 3:6 \\
\cmidrule{3-3}          & 神不是死人的神，乃是活人的神。 & 列王紀上 18:36 \\
    \bottomrule
    \end{tabular}%
  \label{tab:addlabel}%
  \caption{馬太福音引用舊約表格 3} 
\end{table}%


% Table generated by Excel2LaTeX from sheet 't4'
\begin{table}[htbp]
  \centering

    \begin{tabular}{|p{4.2em}|p{31.645em}|p{7.3em}|}
    \toprule
    \rowcolor[rgb]{ .749,  .749,  .749} \textcolor[rgb]{ .267,  .447,  .769}{\textbf{馬太福音}} 
    & \textcolor[rgb]{ .267,  .447,  .769}{\textbf{馬太引用經文內容}} 
    & \textcolor[rgb]{ .267,  .447,  .769}{\textbf{引用出處}} \\
    \midrule
    22:37 & 你要尽心、尽性、尽意爱主你的神。 & 申命記 6:5,10:12 \\
    \midrule
    22:44 & ‘主对我主说：你坐在我的右边，等我把你仇敌放在你的脚下’？ & 詩篇 110:1 \\
    \midrule
    23:17 & 你们这无知瞎眼的人哪！什么是大的？是金子呢，还是叫金子成圣的殿呢？ & 出埃及記 29:37\\
    \multirow{2}[3]{*}{23:19} & 你们这瞎眼的人哪！什么是大的？是礼物呢， & 30:29   \\
\cmidrule{3-3}    \multicolumn{1}{|l|}{} & 还是叫礼物成圣的坛呢？ & 哈該書 2:12 \\
    \midrule
    23:21 & 人指着殿起誓，就是指着殿和那住在殿里的起誓； & 詩篇 132:14 \\
    \midrule
    23:22 & 人指着天起誓，就是指着神的宝座和那坐在上面的起誓。 & 以賽亞書 66:1 \\
    \midrule
    23:24 & 你们这瞎眼领路的，蠓虫你们就滤出来，骆驼你们倒吞下去！ & 以賽亞書 9:16 \\
    \midrule
    \multirow{2}[4]{*}{23:35} & 叫世上所流義人的血都歸到你們身上，從義人亞伯的血起， & 創世記 4:8 \\
\cmidrule{3-3}    \multicolumn{1}{|l|}{} & 直到你們在殿和壇中間所殺的巴拉加的兒子撒迦利亞的血為止。 & 歷代志下 24:21 \\
    \midrule
    23:38 & 看哪，你們的家成為荒場留給你們！ & 耶利米書 22:5 \\
    \midrule
    23:39 & 我告訴你們：從今以後，你們不得再見我，直等到你們說：『奉主名來的是應當稱頌的！』」 & \multicolumn{1}{l|}{詩篇 118:26} \\
    \midrule
    24:15 & 「你們看見先知但以理所說的『那行毀壞可憎的』站在聖地—讀這經的人須要會意— & 但以理書 9:27,11:31,12:11 \\
    \midrule
    \multirow{2}[4]{*}{24:28} & \multirow{2}[4]{*}{屍首在哪裡，鷹也必聚在那裡。} & 以西結書 39:17 \\
\cmidrule{3-3}    \multicolumn{1}{|l|}{} & \multicolumn{1}{l|}{} & 哈巴谷書 1:8 \\
    \midrule
    \multirow{5}[10]{*}{24:29} & \multicolumn{1}{r|}{} & 以賽亞書 13:10,24:23,34:4 \\
\cmidrule{3-3}    \multicolumn{1}{|l|}{} & 「那些日子的災難一過去，『日頭就變黑了，月亮也不放光， & 以西結書 32:7 \\
\cmidrule{3-3}    \multicolumn{1}{|l|}{} & 眾星要從天上墜落，天勢都要震動。』 & 約珥書 2:10,30,31,3:15 \\
\cmidrule{3-3}    \multicolumn{1}{|l|}{} & \multicolumn{1}{r|}{} & 阿摩司書 5:20,8:9, \\
\cmidrule{3-3}    \multicolumn{1}{|l|}{} & \multicolumn{1}{r|}{} & 西番雅書 1:15 \\
    \midrule
    24:30 & 那時，人子的兆頭要顯在天上，地上的萬族都要哀哭。他們要看見人子有能力、有大榮耀駕著天上的雲降臨。 & 但以理書 7:13 \\
    \midrule
    \multirow{2}[4]{*}{24:31} & 當號筒發出響聲，他要差派使者，把他的選民從四方 & 以賽亞書 27:13 \\
\cmidrule{3-3}    \multicolumn{1}{|l|}{} & 從天這邊到天那邊都招聚來。 & 撒迦利亞書 2:6 \\
    \midrule
    24:37-39 & 挪亞的日子怎樣，人子降臨也要怎樣。 & 創世記 7:6 \\
    \bottomrule
    \end{tabular}%
  \label{tab:addlabel}%
    \caption{馬太福音引用舊約表格 4} 
\end{table}%

% Table generated by Excel2LaTeX from sheet 'Sheet8'
\begin{table}[htbp]
  \centering
  \begin{tabular}{|p{4.2em}|p{31.645em}|p{7.3em}|}
    \toprule
    \rowcolor[rgb]{ .749,  .749,  .749} \textcolor[rgb]{ .267,  .447,  .769}{\textbf{馬太福音}} 
    & \textcolor[rgb]{ .267,  .447,  .769}{\textbf{馬太引用經文內容}} 
    & \textcolor[rgb]{ .267,  .447,  .769}{\textbf{引用出處}} \\
  
    \midrule
    26:17 & 除酵節的第一天，門徒來問耶穌說：「你吃逾越節的筵席，要我們在哪裡給你預備？」 & 出埃及記 12:18 \\
    \midrule
    26:28 & 因為這是我立約的血，為多人流出來，使罪得赦。 & 出埃及記 24:8 \\
    \midrule
    26:31 & 那時，耶穌對他們說：「今夜你們為我的緣故都要跌倒，因為經上記著說：『我要擊打牧人，羊就分散了。』 & 撒迦利亞書 13:7 \\
    \midrule
    26:52 & 耶穌對他說：「收刀入鞘吧！凡動刀的，必死在刀下。 & 出埃及記 21:12 \\
    \midrule
    26:60 & 雖有好些人來作假見證，總得不著實據。末後，有兩個人前來說： 61「這個人曾說：『我能拆毀神的殿，三日內又建造起來。』」 & 申命記 19:15 \\
    \midrule
    26:63 & 耶穌卻不言語。大祭司對他說：「我指著永生神叫你起誓告訴我們，你是神的兒子基督不是？」 & 以賽亞書 53:7 \\
    \midrule
    27:9  & 所以那塊田直到今日還叫做「血田」。 9這就應了先知耶利米的話說：「他們用那三十塊錢，就是被估定之人的價錢，是以色列人中所估定的， 10買了窯戶的一塊田，這是照著主所吩咐我的。」 & 撒迦利亞書 11:12,13 \\
    \bottomrule
    \end{tabular}%
  \label{tab:addlabel}%
 \caption{馬太福音引用舊約表格 5}  
\end{table}%





% Table generated by Excel2LaTeX from sheet 'Sheet6'
\begin{table}[htbp]
  \centering

    \begin{tabular}{|p{6.93em}|p{13.93em}|r|r|}
    \toprule
    \rowcolor[rgb]{ .749,  .749,  .749} \textcolor[rgb]{ .267,  .447,  .769}{\textbf{馬太福音}} & \textcolor[rgb]{ .267,  .447,  .769}{\textbf{引用出處}} & \multicolumn{1}{p{6.715em}|}{\textcolor[rgb]{ .267,  .447,  .769}{\textbf{馬太福音}}} & \multicolumn{1}{p{12.645em}|}{\textcolor[rgb]{ .267,  .447,  .769}{\textbf{引用出處}}} \\
    \midrule
    24:29 & 阿摩司書 5:20,8:9 & \multicolumn{1}{p{6.715em}|}{24:28} & \multicolumn{1}{p{12.645em}|}{哈巴谷書 1:8} \\
    \midrule
    12:2  & 出埃及記20：8-11 & \multicolumn{1}{p{6.715em}|}{23:19} & \multicolumn{1}{p{12.645em}|}{哈該書 2:12} \\
    \midrule
    2:15  & 出埃及記 4:22 & \multicolumn{1}{p{6.715em}|}{9:13} & \multicolumn{1}{p{12.645em}|}{何西阿書 6:6} \\
    \midrule
    23:17 & 出埃及記 30:29 & \multicolumn{1}{p{6.715em}|}{2:15} & \multicolumn{1}{p{12.645em}|}{何西阿書 11:1} \\
    \midrule
    17:24 & 出埃及記 30:13, 38:26 & \multicolumn{1}{p{6.715em}|}{12:5} & \multicolumn{1}{p{12.645em}|}{利未記 24:8} \\
    \midrule
    22:32 & 出埃及記 3:6 & \multicolumn{1}{p{6.715em}|}{12:3,4} & \multicolumn{1}{p{12.645em}|}{利未記 24:9} \\
    \midrule
    23:19 & 出埃及記 29:37 & \multicolumn{1}{l|}{5:38} & \multicolumn{1}{p{12.645em}|}{利未記 24:20} \\
    \midrule
    26:28 & 出埃及記 24:8 & \multicolumn{1}{l|}{5:43} & \multicolumn{1}{p{12.645em}|}{利未記 19:18} \\
    \midrule
    \multicolumn{1}{|l|}{5:38} & 出埃及記 21:24 & \multicolumn{1}{p{6.715em}|}{5:27} & \multicolumn{1}{p{12.645em}|}{利未記 18:20} \\
    \midrule
    15:4  & 出埃及記 21:17 & \multicolumn{1}{p{6.715em}|}{8:4} & \multicolumn{1}{l|}{利未記 14:3,4} \\
    \midrule
    26:52 & 出埃及記 21:12 & \multicolumn{1}{p{6.715em}|}{22:32} & \multicolumn{1}{p{12.645em}|}{列王紀上 18:36} \\
    \midrule
    19:19 & 出埃及記 20:16 & \multicolumn{1}{p{6.715em}|}{12:42} & \multicolumn{1}{p{12.645em}|}{列王紀上 10:1} \\
    \midrule
    5:27  & 出埃及記 20:14 & \multicolumn{1}{p{6.715em}|}{2:6} & \multicolumn{1}{p{12.645em}|}{彌迦書 5:2} \\
    \midrule
    5:21  & 出埃及記 20:13 & \multicolumn{1}{p{6.715em}|}{11:14} & \multicolumn{1}{p{12.645em}|}{瑪拉基書 3:1, 4:5} \\
    \midrule
    26:17 & 出埃及記 12:18 & \multicolumn{1}{p{6.715em}|}{11:10} & \multicolumn{1}{p{12.645em}|}{瑪拉基書 3:1} \\
    \midrule
    11:24 & 創世記19：1-29 & \multicolumn{1}{p{6.715em}|}{19:5} & \multicolumn{1}{p{12.645em}|}{瑪拉基書 2:15} \\
    \midrule
    24:37-39 & 創世記 7:6 & \multicolumn{1}{l|}{5:33} & \multicolumn{1}{p{12.645em}|}{民數記 30:2} \\
    \midrule
    23:35 & 創世記 4:8 & \multicolumn{1}{p{6.715em}|}{2:15} & \multicolumn{1}{p{12.645em}|}{民數記 24:8} \\
    \midrule
    19:4  & 創世記 1:27,5:2 & \multicolumn{1}{p{6.715em}|}{21:5} & \multicolumn{1}{p{12.645em}|}{撒迦利亞書 9:9} \\
    \midrule
    1:1-6 & 代上 2:1-14，3:5-24 & \multicolumn{1}{p{6.715em}|}{24:31} & \multicolumn{1}{p{12.645em}|}{撒迦利亞書 2:6} \\
    \midrule
    12:42 & 歷代志下 9:1 & \multicolumn{1}{p{6.715em}|}{26:31} & \multicolumn{1}{p{12.645em}|}{撒迦利亞書 13:7} \\
    \midrule
    19:5  & 歷代志下 9:1 & \multicolumn{1}{p{6.715em}|}{27:9} & \multicolumn{1}{p{12.645em}|}{撒迦利亞書 11:12,13} \\
    \midrule
    23:35 & 歷代志下 24:21 & \multicolumn{1}{p{6.715em}|}{12:3,4} & \multicolumn{1}{p{12.645em}|}{撒母耳記上 21:4,6} \\
    \midrule
    24:15 & 但以理書 9:27,11:31,12:11 &       &  \\
    \midrule
    24:30 & 但以理書 7:13 &       &  \\
    \bottomrule
    \end{tabular}%
  \label{tab:addlabel}%
   \caption{馬太福音引用舊約表格 6} 
\end{table}%

% Table generated by Excel2LaTeX from sheet 'sorted2'
\begin{table}[htbp]
  \centering

    \begin{tabular}{|p{7.215em}|p{10.645em}|p{8.645em}|p{12.285em}|}
    \toprule
    \rowcolor[rgb]{ .749,  .749,  .749} \textcolor[rgb]{ .267,  .447,  .769}{\textbf{馬太福音}} & \textcolor[rgb]{ .267,  .447,  .769}{\textbf{引用出處}} & \textcolor[rgb]{ .267,  .447,  .769}{\textbf{馬太福音}} & \textcolor[rgb]{ .267,  .447,  .769}{\textbf{引用出處}} \\
    \midrule
    4:4   & 申命記 8:3 & 21:13 & 耶利米書 7:11 \\
    \midrule
    22:37 & 申命記 6:5,10:12 & 2:18  & 耶利米書 31:15 \\
    \midrule
    4:7   & 申命記 6:16 & 23:38 & 耶利米書 22:5 \\
    \midrule
    4:10  & 申命記 6:13， 10:20 & 1:22  & 以賽亞書7：14 \\
    \midrule
    19:19 & 申命記 5:20 & 23:24 & 以賽亞書 9:16 \\
    \midrule
    19:19 & 申命記 5:19 & 4:15  & 以賽亞書 9:1 \\
    \midrule
    5:27  & 申命記 5:18 & 23:22 & 以賽亞書 66:1 \\
    \midrule
    5:21  & 申命記 5:17 & 21:5  & 以賽亞書 62:11 \\
    \midrule
    5:31  & 申命記 24:1，3 & 13:15 & 以賽亞書 6:10 \\
    \midrule
    19:7  & 申命記 24:1 & 21:13 & 以賽亞書 56:7 \\
    \midrule
    \multicolumn{1}{|l|}{5:33} & 申命記 23:23 & 26:63 & 以賽亞書 53:7 \\
    \midrule
    \multicolumn{1}{|l|}{5:38} & 申命記 19:21 & 8:17  & 以賽亞書 53:4 \\
    \midrule
    26:60 & 申命記 19:15 & 12:18-21 & 以賽亞書 42:1-4 \\
    \midrule
    4:6   & 詩篇 91:11，12 & 3:3   & 以賽亞書 40:3 \\
    \midrule
    21:16 & 詩篇 8:2 & 11:5  & 以賽亞書 35：5 \\
    \midrule
    13:35 & 詩篇 78:2 & 15:8,9 & 以賽亞書 29:13 \\
    \midrule
    23:21 & 詩篇 132:14 & 21:42 & 以賽亞書 28:16 \\
    \midrule
    21:9  & 詩篇 118:26 & 24:31 & 以賽亞書 27:13 \\
    \midrule
    23:39 & \multicolumn{1}{l|}{詩篇 118:26} & 24:29 & 以賽亞書 13:10,24:23,34:4 \\
    \midrule
    21:42 & 詩篇 118:22 & 24:28 & 以西結書 39:17 \\
    \midrule
    22:44 & 詩篇 110:1 & 24:29 & 以西結書 32:7 \\
    \midrule
    24:29 & 西番雅書 1:15 & 24:29 & 約珥書 2:10,30,31,3:15 \\
    \midrule
          &       & 12:39-40 & 約拿書  \\
    \midrule
          &       & 2:23  & ? \\
    \bottomrule
    \end{tabular}%
  \label{tab:addlabel}%
    \caption{馬太福音引用舊約表格 7} 
\end{table}%
%\begin{wraptable}{r}{5.5cm}
%\caption{A wrapped table going nicely inside the text.}\label{wrap-tab:1}
%\begin{tabular}{ccc}\\\toprule  
%Header-1 & Header-1 & Header-1 \\\midrule
%2 &3 & 5\\  \midrule
%2 &3 & 5\\  \midrule
%2 &3 & 5\\  \bottomrule
%\end{tabular}
%\end{wraptable} 

\subsubsection{馬可福音}

\subsubsection{路加福音}

 \subsubsection{约翰福音}   
    
    



%\begin{tabular}{|l|c|}
%\rowcolor[gray]{.9}
%one&two\\
%\rowcolor[gray]{.5}
%three&four
%\end{tabular}


% Table generated by Excel2LaTeX from sheet 'Sheet1'
%\begin{table}[htbp]
%  \centering
 % \caption{Add caption}
 %   \begin{tabular}{|llll|}
 %   \toprule
 %   \rowcolor[rgb]{ .267,  .447,  .769} \textcolor[rgb]{ 1,  1,  1}{\textbf{34}} & %\textcolor[rgb]{ 1,  1,  1}{\textbf{df}} & \textcolor[rgb]{ 1,  1,  1}{\textbf{rt}} & \textcolor[rgb]{ 1,  1,  1}{\textbf{fg}} \\
%    \midrule
 %   \rowcolor[rgb]{ .851,  .882,  .949} 地方    & 對方溝通  & 能不能   & 方法和 \\
 %   \midrule
  %  下浮郡通過 & 買回家   & 明白你不敢 & 規劃結構 \\
  %  \bottomrule
  %  \end{tabular}%
 % \label{tab:addlabel}%
%\end{table}%


\subsection{兩約之間}
聖經的舊約止於以色列亡國被擄與歸回，當時歸回的主要是猶大，便雅憫，和利未支派，他們被稱為猶太人。從那時起，一直到耶穌的時代，通稱為兩約之間，亦即新約與舊約之間。

\subsubsection{歸回時期：波斯統治時代，公元前536-332年}
波斯王古列時期，猶太人第一次歸回，由大衛的後裔所羅巴伯與祭司約書亞領導。他們一回到猶大地，立刻著手重建聖殿，但是才開工沒多久，就遇到鄰邦阻撓，建殿聖工遂遭擱置。幸好此時有先知哈該和撒迦利亞及時起來鼓勵百姓，他們才在大利烏王一世第二年復工，而於第六年完工。這段建殿經過記載在舊約的以斯拉記裡。猶太人險遭滅族，末底改和以斯帖挺身而出，整個民族得蒙上帝奇妙拯救。文士以斯拉帶領第二批猶太人歸回，他搜集、考究律法，並在聖殿裡負責教導律法。省長尼希米第三次归回，率領百姓重建耶路撒冷的城牆，並整頓當時一些弊病。最後，瑪拉基提醒百姓歸向上帝，並且預言主的使者即將來到。舊約就寫到瑪拉基書為止。

\subsubsection{熬煉時期：希臘統治時代，公元前332-167年}
希臘在亞歷山大的率領之下，征服波斯，建立一個強大的帝國。但是他在 33 歲時英年早逝，帝國分裂為四，其中兩個與猶太地有密切關係，亦即南邊的埃及多利買王朝與北邊的敘利亞西流古王朝，此兩王朝在文化上均承襲希臘的影響。

首先，猶太在埃及多利買王朝的統管之下(公元前 323-198年)。此時，埃及的統治者仍為希臘人，不過，其推行希臘化的手段比較溫和，而且不干預猶太人的宗教。多利買王朝的非拉鐵非王(公元前 285-246 年)善待猶太人，致力保存猶太文化，甚至花錢買贖猶太奴隸。最重要的是他召集猶太學者完成舊約翻譯工作，舊約希臘文七十士譯本就在非拉鐵非王的授命下完成。

接著，敘利亞的西流古王朝接管猶太(公元前 198-167 年)。原先，猶太人歡迎安提阿古 III 的統治，但後來他與羅馬交戰，戰敗需賠鉅款，遂向猶太人索重稅，而引起猶太人的不滿。他把巴勒斯坦及鄰近地區劃分為加利利、猶太、撒馬利亞、特拉可尼等區。等安提阿古 IV 繼位時(公元前 186 年)，猶太人可以說開始了一段艱苦的歲月。西流古王朝致力於傳播希臘文化。在猶太人中間為了要不要接受希臘化，已經產生不同看法。有些顯赫富有的人，為了迎合統治者，就全盤接受這種新的生活方式，可謂當時之新潮派；他們則自稱為進步黨或自由派。另一些比較保守的人，堅持猶太傳統文化及信仰，他們自稱為哈西典人，意思就是敬虔派。安提阿古 IV 極力推行希臘化，他傲慢又兇殘，向猶太人索重稅，並且兩次接受賄賂，任意更換大祭司，因此遭猶太人抵制。有一次他和埃及王打仗，猶太人聽說他陣亡，大開慶祝會，不料他活著回來，見狀大怒，更加逼迫猶太人，不准他們行割禮、守安息日，並擄掠耶路撒冷。據說他殺了四萬，又俘虜四萬。最叫猶太人難以消受的是，他在耶路撒冷聖殿裡安放宙斯(Zeus)的神像，立宙斯的祭壇，又強迫猶太人獻豬給宙斯，時為公元前 168 年 12 月 25 日。猶太人對他可以說恨之入骨。

\subsubsection{中興時期：馬加比革命時期，公元前167-63年}
安提阿古 IV 不但在耶路撒冷這樣做，還要把這種拜偶像的習俗推行到各地。公元前 167 年，有一個敘利亞軍官率領一批士兵來到耶路撒冷城西的莫頂(Modin)小城，高築希臘神像，並召人民來獻祭。他對當地極有名望的老祭司馬他提亞(Mattathias)威迫利誘，要他帶頭順服國王，獻祭給神像。馬他提亞不從，但有一個猶太人為了討好長官，竟然上前獻祭，馬他提亞見狀，怒不可遏，當場就把他給殺了。他的兒子們一不做，二不休，又把敘利亞軍官也殺了。經他們登高一呼，凡是願意敬拜上帝，遵守律法，不願意聽從國王去拜偶像的猶太人，都來跟隨他們。馬他提亞此時年紀已經老邁，臨終時指定三子猶大繼續領導。猶大驍勇善戰，大家暱稱他為馬加比(Maccabees)，意思就是「執鐵鎚者」，因此這段事蹟就被稱為「馬加比革命」。

兩年後，也就是公元前 165 年的 12 月 25日，猶太人在耶路撒冷潔淨聖殿，去除偶像，再一次將聖殿奉獻給上帝，這就是耶穌時期「修殿節」的由來。猶太人現在稱它為「看努加」，亦即「奉獻節」的意思；又由於過節時家家戶戶點臘燭慶祝，於是又叫做「眾光節」。

之後，猶太人在馬他提亞五個兒子們逐一領導之下，繼續奮鬥﹔雖然敘利亞王屢次派軍隊鎮壓，皆未成功。到馬他提亞其中一個兒子西門任領袖時，他被猶太人擁立為終身大祭司。此時敘利亞國力稍弱，需要猶太人的支持，遂容許猶太人自治。西門在公元前 139 年又與羅馬結盟，羅馬承認猶大國的獨立自主權。但是未多久，西門和他的兩個兒子被他的女婿暗殺，僅存的兒子約翰赫迦納及時登基，赫迦納終成大祭司和國家元首，這就開始了哈斯摩年王朝。

\subsubsection{亡國時期：羅馬統治時期，公元前63─公元后70年}

猶太人在敘利亞西流古王朝沒落之後，經過一百多年光榮的獨立時期。其實猶太人真正的敵人不是外來的侵略者，而是內部的不團結。當時，猶太人有兩大黨派，一個是支持馬加比家族作君王和祭司的人；他們擁有政治權勢，是耶穌時代撒都該人的前身。另一個就是前面所提到的哈西典人。他們原先也支持馬加比家族，但是在爭取到宗教方面的自由之後，即不再過問政治；法利賽人就是從他們演變而來。這兩派的人經常有爭執，而哈斯摩年王朝更是為了王位的繼承，內部迭有紛爭。當時，羅馬正逐漸興起，猶太最後敵不過羅馬的勢力，終為羅馬的龐貝將軍所奪。

此時有以土買人(亦即以東人)安提帕崛起，他被龐貝授命為哈斯摩年王朝的首相和顧問，在猶太地操有實權。但是他同時卻又幫助凱撒，因此當凱撒後來戰勝龐貝成為羅馬皇帝時，就委派安提帕為猶太、撒馬利亞和加利利三省的省長，其兒子希律後來更被封為王。哈斯摩年王朝到此告一段落，其子孫只能擔任大祭司職。希律，就是後來被稱為大希律的，四處逢迎，起初極力援助安東尼 (Antony)，後來又見風轉舵，支持屋大維 (Octavian)，因此，不論誰當羅馬皇帝，希律都得到他們的信任。屋大維，亦即後來被尊稱為奧古斯都的，當上羅馬皇帝以後，更把猶太附近的撒馬利亞、加利利、推羅海濱，甚至遠到埃及，都交給希律管轄，並正式任命他為王， 時為公元前37 年。希律為討好猶太人，在公元前 20 年為他們蓋了一座宏偉的聖殿，但聖殿尚未蓋成，他就在公元前 4 年去世，他的兒子們被封為巴勒斯坦不同地區的分封王。

猶太人雖然有這些外來的分封王，但是真正控制他們，得其心的是祭司。大祭司的職位原是世襲，而且是終身職。與祭司共同治理百姓的還有由長老們組成的議會。除了納貢和外交事務外，猶太人其實享有相當大的自主權。

但是猶太人還是不滿羅馬的干預與統治，陸續有叛變，惟其內部卻又不和；公元70 年羅馬的提多將軍摧毀耶路撒冷，其軍隊焚燒聖殿，猶太人自殺或被殺的很多，猶太史家約瑟夫(Flavius Josephus)所寫的書《猶太戰史》對此有詳細描述。猶太人最後一次起義是在公元 135 年，由巴克巴(Bar-cachba)所領導，導火線是因當時羅馬皇帝哈德良 (Hadrian, 117-138 A.D.)下令禁止行割禮，並在耶路撒冷聖殿舊址建造朱彼特(Jupiter)神廟。但此叛變並未成功，巴克巴被捕。從此以後，耶路撒冷變成羅馬人的城市，猶太至此可以說是徹底覆亡。
　　
\section{宗教狀況}
\subsection{聖殿，會堂}

\subsection{律法}
\subsection{節期}
\subsection{信仰派別}
\subsubsection{法利賽人}

法利賽人的意思是「分別出來」，他們是與當政者不合而分離出來的一班虔誠人士，約在公元前 150 年形成。他們遵行律法，相信死人復活。在耶穌出生前一世紀，法利賽人中有兩位極具影響力的大拉比 ，形成法利賽的兩大學派，由希列(Hillel) 和沙買(Shammai)所領導。希列生於巴比倫，後來遷居耶路撒冷，對律法採取比較自由的態度。沙買這一派比較嚴格，他們反對那些傾向羅馬的政黨，同情並贊助謀求國家獨立。使徒保羅的師傅迦瑪列乃是希列的孫兒。

馬太福音 
23:1-7 那時，耶穌對眾人和門徒講論，說：「文士和法利賽人坐在摩西的位上，3凡他們所吩咐你們的，你們都要謹守遵行，但不要效法他們的行為，因為他們能說不能行。4他們把難擔的重擔捆起來，擱在人的肩上，但自己一個指頭也不肯動。 5他們一切所做的事都是要叫人看見，所以將佩戴的經文做寬了，衣裳的穗子做長了；6喜愛筵席上的首座、會堂裡的高位， 7又喜愛人在街市上問他安，稱呼他『拉比』。
路加福音 20：45-47
45 眾百姓聽的時候，耶穌對門徒說： 46 「你們要防備文士。他們好穿長衣遊行，喜愛人在街市上問他們安，又喜愛會堂裡的高位、筵席上的首座。 47 他們侵吞寡婦的家產，假意作很長的禱告。這些人要受更重的刑罰！」

馬可福音  
醫治癱子2：1-12 有幾個文士坐在那裡，心裡議論說： 7「這個人為什麼這樣說呢？他說僭妄的話了！除了神以外，誰能赦罪呢？」
馬太9:1~8；路加5:18~26

路加福音
辯駁耶穌的權柄 20：1-8  有一天，耶穌在殿裡教訓百姓、講福音的時候，祭司長和文士並長老上前來，
撒都該人辯駁復活之事 20：27-40  有幾個文士說：「夫子，你說得好！」 
凶惡園戶的比喻，納稅給愷撒  19 文士和祭司長看出這比喻是指著他們說的，當時就想要下手拿他，只是懼怕百姓。。。 22 我們納稅給愷撒可以不可以？」

\subsubsection{撒督該人}
撒都該人可能源自所羅門王時代的撒都，他們多為祭司階層的人，可以說是當時的執政者與權貴。他們崇尚理性，只接受摩西五經，不相信死人復活，反對法利賽人在律法之外又加上許多遺傳。其實撒都該人的興趣在於政治，多過宗教。

\subsubsection{奮銳當人}
奮銳黨主張與羅馬對抗，不惜使用武力。但到公元第二世紀猶太覆亡之後，大部份的黨派都消失無蹤

\subsubsection{隱士派}
\subsubsection{希律黨}
希律黨是擁護希律做王的人

\section{百姓生活}

\subsection{希臘文化的影響}
\subsection{職業}


\subsubsection{拉比}
就是老師的意思。
\subsubsection{漁夫}
\subsubsection{稅吏}
\subsubsection{木匠}


\subsection{身體生病的}
\subsection{靈裡受捆綁的}
\subsection{道德生活}

\subsection{生活習俗}
\subsubsection{婚姻}
他母親馬利亞已經許配了約瑟，還沒有迎娶(馬太1：18）
婚宴喝酒
\subsubsection{葬禮}
死人埋石頭洞裡，裹屍布，沒藥
\subsubsection{貨幣：}

\subsubsection{銀行}


\section{政治背景}
\subsection{羅馬政府}
\subsubsection{彼拉多}
\subsubsection{羅馬兵丁，百夫長}

\subsection{分封王希律}
大希律靠其政治手腕被羅馬皇帝委任為猶太人的王，他曾經要殺嬰孩耶穌，沒有成功。他死了以後，遺囑將王國分給他三個兒子，見圖一。當然，這還需經過奧古斯都皇帝的批准。

亞基老	管理猶太、撒馬利亞和以土買等地，他不能稱為王，只能稱為總督(ethnarch)。
希律安提帕	管理加利利和比利亞，他的地位是分封王(tetrarch)，權力比總督小
腓力 I	管理巴勒斯坦北部及東北部，包括特拉可尼，以土利亞等地。

\subsection{宗教領袖}
\subsubsection{祭司}

祭司的職責原是獻祭、傳達上帝的旨意並教誨民眾。以色列人從巴比倫歸回以後，沒有王，外邦人也不希望大衛王室抬頭，祭司的權力便逐漸擴張而擁有政權。馬加比革命之後所建立的王朝更是祭司主政。但是在政治環境爭權奪利的風氣之下，祭司也難保清高。而外來的統治者有時任意廢大祭司，另立新祭司，甚至還有人用錢收買而得大祭司職。
祭司則多半為撒都該人，亦有少數為法利賽人
\subsubsection{文士}

文士原是抄寫律法的人，由利未人擔任，他們保存聖經，不遺餘力。以色列人被擄分散在各處設立會堂時，文士研讀並解釋律法，在教導百姓遵行律法上扮演極重要角色。像以斯拉就是文士，又是祭司。但是後來祭司慢慢忽略律法，文士就與祭司分開，而形成另一個專業。
文士多半為法利賽人，但亦有少數為撒都該人

馬太福音 
23:1-7 那時，耶穌對眾人和門徒講論，說：「文士和法利賽人坐在摩西的位上，3凡他們所吩咐你們的，你們都要謹守遵行，但不要效法他們的行為，因為他們能說不能行。4他們把難擔的重擔捆起來，擱在人的肩上，但自己一個指頭也不肯動。 5他們一切所做的事都是要叫人看見，所以將佩戴的經文做寬了，衣裳的穗子做長了；6喜愛筵席上的首座、會堂裡的高位， 7又喜愛人在街市上問他安，稱呼他『拉比』。

\subsubsection{公會}
詩篇 111:1 你們要讚美耶和華！我要在正直人的大會中，並公會中，一心稱謝耶和華

馬太福音 
10:17 你們要防備人，因為他們要把你們交給公會，也要在會堂裡鞭打你們，
26:59 祭司長和全公會尋找假見證控告耶穌，要治死他。

馬可福音 14:55 祭司長和全公會尋找見證控告耶穌，要治死他，卻尋不著

路加福音 22:66 天一亮，民間的眾長老連祭司長帶文士都聚會，把耶穌帶到他們的公會裡，

約翰福音 11:47 祭司長和法利賽人聚集公會，說：「這人行好些神蹟，我們怎麼辦呢？

使徒行傳 
4:15 於是吩咐他們從公會出去，就彼此商議說：
5：21，27，33，34 使徒听了这话，天将亮的时候就进殿里去教训人。大祭司和他的同人来了，叫齐公会的人和以色列族的众长老，就差人到监里去，要把使徒提出来。 21 他們聽從這話，在清...司問他們說： 27 带到了，便叫使徒站在公会前。大祭司问他们说： 27 他們把使徒們帶來後，叫他們站在議會當中。大祭司質問他們， 27 他们把使徒们带来后，叫他...33 公會的人聽見就極其惱怒，想要殺他們。 33 公会的人听见就极其恼怒，想要杀他们。 33 議會的人聽了，勃然大怒. 34 但有一個法利賽人名叫迦瑪列，是眾百姓所敬重的教法師，在公會中站起來，吩咐人把使徒暫且帶到外面去，
6:12，15 12 他們又聳動了百姓、長老並文士，就忽然來捉拿他，把他帶到公會去  15 在公會裡坐著的人都定睛看他，見他的面貌好像天使的面貌
23：1，6，15，20： 保罗定睛看着公会的人，说：“弟兄们，我在...看出大众一半是撒都该人，一半是法利赛人，就在公会中大声说：“弟兄们，我是法利赛人，也是法利赛人的子孙。我现在受审问，是为盼望死人复活。” 6 保羅知道有一部分人是撒都該人，另一...你們和公會要知會千夫長，叫他帶下保羅到你們這裡來，假作要詳細察考他的事。我們已經預備好了，不等他來到跟前就殺他。」 15 现在你们和公会要知会千夫长，叫他带下保罗到你们这里来  20 他說：「猶太人已經約定，要求你明天帶下保羅到公會裡去，假作要詳細查問他的事。
24:20 即或不然，這些人若看出我站在公會前有妄為的地方，他們自己也可以說明。

\section{地理狀況}

猶太地(Judaea)是希臘和羅馬對猶大(Judah)的稱呼。猶大原為以色列分裂後的南國，亡國被擄之後被併入撒馬利亞省，尼希米時期則單獨成一省，有自己的省長。

希臘羅馬時期的猶太(Judaea)廣義來說，它涵蓋整個巴勒斯坦，包括加利利和撒馬利亞；狹義的意思則只包括原先猶大的地區。猶太地(Judaea)和信奉猶太宗教(Judaism)的猶太人(Jews)，三者中文易混淆，需加區分。


\section{疑難問題}
\subsection{馬太福音}
1：1-29， 家譜·:為什麼家譜用14代，為什麼省略了幾個王？
27：9，所以那塊田直到今日還叫做「血田」。 這就應了先知耶利米的話說：「他們用那三十塊錢，就是被估定之人的價錢，是以色列人中所估定的， 買了窯戶的一塊田，這是照著主所吩咐我的。」， 是撒迦利亞書裡的話，為什麼說是耶利米書呢？

\subsection{馬可福音}
\subsection{路加福音}
\subsection{約翰福音}
\subsubsection{洗澡与洗脚的区别 13}
\begin{itemize}
\item 旧约

\begin{itemize}
\item 立亞倫及其子為祭司當行之禮 出29：1-7 \\
你使亞倫和他兒子成聖，給我供祭司的職分，要如此行。…. 要使亞倫和他兒子到會幕門口來，用水洗身。 …將聖冠加在冠冕上，就把膏油倒在他頭上膏他。
\item 每天為百姓服務時需來一次洗手洗脚潔淨禮 出30:18-21\\
耶和華曉諭摩西說：「你要用銅做洗濯盆和盆座，以便洗濯。要將盆放在會幕和壇的中間，在盆裡盛水。亞倫和他的兒子要在這盆裡洗手洗腳。他們進會幕，或是就近壇前供職給耶和華獻火祭的時候，必用水洗濯，免得死亡。他們洗手洗腳，就免得死亡。這要做亞倫和他後裔世世代代永遠的定例。」
\end{itemize}

\item 新约
\begin{itemize}
\item 洗澡是洗身比喻得救；這次的經歷不用重複一次足可。當人信主時，他的罪全被塗抹洗淨。
\begin{itemize}
\item 林前6:9-11\\
你們豈不知不義的人不能承受神的國嗎？….11 你們中間也有人從前是這樣，但如今你們奉主耶穌基督的名，並藉著我們神的靈，已經洗淨、成聖、稱義了。
\item 提多 3:4-7\\
到了神我們救主的恩慈和他向人所施的慈愛顯明的時候， 5 他便救了我們，並不是因我們自己所行的義，乃是照他的憐憫，藉著重生的洗和聖靈的更新。…好叫我們因他的恩得稱為義，可以憑著永生的盼望成為後嗣。
\item 啟示錄 1:5-6\\
他愛我們，用自己的血使我們脫離(洗去)罪惡， 6 又使我們成為國民，做他父神的祭司。
\end{itemize}

\item 洗腳是洗身體的部份，比喻生活或事奉。但是他的生活常與世界接觸，甚易被世界沾染，以致屢常犯罪，所以要常求神潔淨、赦免，這是腳的部份，因經常走路，被塵埃泥濘弄污，故要經常清洗。
\begin{itemize}
\item 約壹1:9\\
我們若認自己的罪，神是信實的，是公義的，必要赦免我們的罪，洗淨我們一切的不義。
\end{itemize}
\end{itemize}
\end{itemize}


